\chapter{Authentifikation}

Wenn sich Alice in einen Hostrechner einloggt, wie kann er
Rechner wissen, da\3 es wirklich Alice ist? Wie kann er
erkennen, da"s Eve versucht sich als Alice auszugeben?
"Ublicherweise wird dieses Problem mit Hilfe eines Passworts
gel"ost. Alice gibt ihr Passwort ein und der Rechner "uberpr"uft, ob es
korrekt ist. Alice und der Rechner haben ein gemeinsames Geheimnis, das
der er jedesmal von Alice verlangt, wenn sie sich einloggen will.

%%
%% Authentifikation mit Einwegfunktionen
%%

\section{Authentifikation mit Einwegfunktionen}

Roger Needham und Mike Guy bemerkten, da\3 der Host gar nicht die
Passw"orter zu kennen braucht. Der Host mu\3 nur die richtigen
Passw"orter von den falschen unterscheiden k"onnen. Dies kann einfach
mit Einwegfunktionen realisiert werden. Der Host speichert nur die
Werte $f(Passwort)$, wobei $f$ eine Einwegfunktion ist.
\begin{enumerate}
  \item Alice sendet ihr Passwort dem Host.
  \item Der Host wendet eine Einwegfunktion auf das Passwort an.
  \item Der Host vergleicht das Ergebnis mit dem zuvor gespeicherten Wert.
\end{enumerate}
Da jetzt der Host nicht die Passw"orter speichert, sondern nur die
Werte der mit einer Einwegfunktion verschl"usselten Passw"orter, ist
diese Liste f"ur einen Angreifer nicht direkt verwendbar, da er daraus die
Passw"orter nicht mehr so leicht rekonstruieren kann.

Diese Authentifikation ist nicht sonderlich sicher, da ein {\em
W"orterbuchangriff} unternommen werden kann.  Eine Datei, die mit
einer Einwegfunktion verschl"usselten Passw"orter enth"alt, stellt
immer noch einen Schwachpunkt dar. Mallory erstellt eine Liste der am
h"aufigsten benutzten Passw"orter, verschl"usselt jedes Passwort mit
der Einwegfunktion und speichert die Ergebnisse. Mallory schafft es die
Datei mit den verschl"usselten Passw"ortern, die sich auf dem
Hostrechner befindet, zu stehlen und sucht nach "Ubereinstimmungen in
den beiden Dateien. Dieser Angriff ist erstaunlich erfolgreich, was
nicht an dem Verfahren mit der Einwegfunktion liegt, sondern ist den
unachtsamen Benutzern anzulasten, da sie einfach erratbare Passw"orter
w"ahlen.

%%
%% Authentifikation mit Public Key
%%

\section{Authentifikation mit Hilfe von Public Key}

Die Protokolle zur Authentifikation mit Hilfe von Passw"ortern haben
eine weitere enorme Sicherheitsl"ucke. Wenn Alice das Passwort zu
ihrem Host sendet, kann Eve ihr Passwort lesen, wenn sie Zugang zu dem
Kommunikationsweg oder dem Prozessorspeicher des Hosts hat.

Mit Hilfe der Public-Key Kryptographie kann dieses Problem behoben
werden. Der Host hat eine Datei mit den "offentlichen Schl"usseln
aller Benutzer und alle Benutzer haben ihren "offentlichen
Schl"ussel. Ein einfaches Protokoll kann wie folgt aussehen:
\begin{enumerate}
  \item Der Host sendet Alice eine Zufallszahl $r$.
  \item Alice verschl"usselt $r$ mit ihrem geheimen Schl"ussel und
        schickt das Chiffrat an den Host.
  \item Der Host entschl"usselt das Chiffrat mit dem "offentlichen
        Sch"ussel von Alice und vergleicht es mit $r$.
\end{enumerate}
Ein Angreifer kann sich jetzt nicht dem Host gegen"uber als Alice
identifizieren, da er keinen Zugang zu Alices geheimen Schl"ussel
hat. Alice sendet niemals ihren geheimen Schl"ussel "uber den
Kommunikationsweg zu dem Host. Eve, die die Interaktion zwischen Alice
und dem Host abh"ort, kann keine Information gewinnen, die es ihr
erm"oglichen w"urde, Alices geheimen Schl"ussel zu berechnen oder sich
als Alice auszugeben, ohne da\3 es der Host bemerkt. Weder der Host
noch der Kommunikationsweg m"ussen sicher sein.

Es ist f"ur Alice nicht geschickt, eine beliebige Zufallszahl $r$, die
sie zugeschickt bekommt, zu verschl"usseln. Eve k"onnte dadurch
die an Alice geschickten Nachrichten entschl"usseln. Ein verbessertes 
Protokoll sieht wie folgt aus:
\begin{enumerate}
  \item Alice berechnet etwas, wobei sie eine eigene Zufallszahl und
        ihren geheimen Schl"ussel verwendet, und schickt das Ergebnis
        an den Host.
  \item Der Host sendet Alice eine andere Zufallszahl.
  \item Alice verwendet eine vorgegebene Funktion, um einen
        Funktionswert, der von den beiden Zufallszahlen und ihrem
        geheimen Schl"ussel abh"angt, zu berechnen und schickt das
        Ergebnis an den Host.
  \item Der Host "uberpr"uft dies.
\end{enumerate}
Falls Alice genauso wenig dem Host vertraut, wie der Host Alice traut,
fordet sie ihn auf, sich ihr gegen"uber zu authentifizieren.

%%
%% Authentifikation mit Schl"usselaustausch
%%

\section{Authentifikation mit Schl"usselaustausch}

Diese Protokolle verbinden die Authentifikation mit einem
Schl"usselaustausch, um ein generelles Problem zu l"osen: Alice und
Bob sind an zwei verschiedenen Enden eines Netzwerks und wollen sicher
Daten "ubertragen. Wie k"onnen Alice und Bob einen geheimen Schl"ussel
austauschen und dabei sicher sein, da\3 keiner von ihnen in
Wirklichkeit mit einem Angreifer Mallory spricht? Die meisten dieser
Protokolle basieren auf der Tatsache, da\3 eine Schl"usselzentrale mit
jedem Teilnehmer einen gemeinsamen geheimen Schl"ussel besitzt und
da\3 alle Schl"ussel am richtigen Ort sind, bevor das Protokoll
beginnt.

%%
%% Wide-mouth frog
%%

\subsubsection*{Wide-mouth frog}

Das wide-mouth frog Protokoll ist vielleicht das einfachste Protokoll,
das ein symmetrisches Verfahren und eine vertrauensw"urdige
Schl"usselzentrale benutzt.  Alice und Bob besitzen jeweils einen
geheimen Schl"ussel $A$ bzw. $B$ mit der Schl"usselzentrale. Diese
Schl"ussel werden nur f"ur den Schl"usselaustausch und nicht zur
Verschl"usselung von eigentlichen Nachrichten verwendet. Mit Hilfe von
nur zwei Nachrichten "ubermittelt Alice an Bob einen geheimen
Sitzungsschl"ussel $K$.  Mit $T_A$ und $T_B$ werden Zeitstempel
bezeichnet. Die Verschl"usselungsfunktionen $E_A$ und $E_B$ sind mit
den geheimen Schl"usseln parametrisiert, die Alice und Bob mit der
Schl"usselzentrale teilen.
\begin{enumerate}
  \item Alice erzeugt das Tupel $(Alice,E_A(T_A,Bob,K))$ und schickt es an
        die Schl"us\-sel\-zen\-trale.
  \item Die Schl"usselzentrale entschl"usselt die Nachricht von
        Alice. Sie erzeugt einen neuen Zeitstempel $T_B$ und sendet das
        verschl"usselte Tupel $E_B(T_B,Alice,K)$ an Bob.
\end{enumerate}
Die wichtigste Annahme bei diesem Protokoll ist, da\3 Alice kompetent
ist, gute (sichere) Sitzungschl"ussel zu erzeugen. Man mu\3 beachten,
da\3 es schwierig ist, Zufallszahlen zu erzeugen. Alice ist
m"oglicherweise nicht dazu in der Lage.

%%
%% Needham-Schroeder
%%

\subsubsection*{Needham-Schroeder}

Das Protokoll von Needham und Schroeder benutzt ein symmetrisches Verfahren
und eine Schl"usselzentrale.
\begin{enumerate}
  \item Alice erzeugt eine Zufallszahl $R_A$ und schickt das Tupel
        $(Alice,Bob,R_A)$ an die Schl"us\-sel\-zen\-trale.
  \item Die Schl"usselzentrale erzeugt einen zuf"alligen
        Sitzungsschl"ussel $K$ und schickt das Tupel
        $E_A(R_A,Bob,K,E_B(K,Alice))$ an Alice.
  \item Alice entschl"usselt und erh"alt den Sitzungsschl"ussel
        $K$. Sie "uberpr"uft die Zufallszahl $R_A$ und schickt das
        Tupel $E_B(K,Alice)$ an Bob.
  \item Bob entschl"usselt und erh"alt den Sitzungsschl"ussel $K$. Er
        erzeugt eine Zufallszahl $R_B$ und schickt das Tupel
        $E_K(R_B)$ an Alice.
  \item Alice entschl"usselt und erh"alt die Zufallszahl $R_B$. Sie
        berechnet $R_B-1$ und schickt das Tupel $E_K(R_B-1)$ an Bob.
  \item Bob pr"uft, ob $R_B-1$ geschickt wurde.
\end{enumerate}

%%
%% Kerberos
%%

\subsubsection*{Kerberos}

Das Kerberos-Protokoll ist eine Variante des
Needham-Schroeder-Protokolls. Die Zufallszahl $R_A$ wird durch einen
Zeitstempel ersetzt.
\begin{enumerate}
  \item Alice sendet das Tupel $(Alice,Bob)$ an die
        Schl"usselzentrale.
  \item Die Schl"usselzentrale schickt das Tupel
        $(E_A(T_{SZ},L,K,Bob),E_B(T_{SZ},L,K,Alice))$ an Alice. Die
        Nachricht wird mit einer G"ultigkeitsdauer $L$ und einem
        Zeitstempel $T_{SZ}$ versehen.
  \item Alice schick das Tupel
        $(E_K(Alice,T_A),E_B(T_{SZ},L,K,Alice))$ an Bob.
  \item Bob sendet das Tupel $E_K(T_A+1)$ an Alice.
\end{enumerate}

%%
%% Neuman-Stubblebine
%%

\subsubsection*{Neuman-Stubblebine}

\begin{enumerate}
  \item Alice erzeugt eine Zufallszahl $R_A$ und sendet das Tupel
        $(Alice,R_A)$ an Bob.
  \item Bob erzeugt eine Zufallszahl $R_B$ und sendet das Tupel
        $(Bob,R_B,E_B(Alice,R_A,T_B))$ mit einem Zeitstempel $T_B$
        versehen an die Schl"usselzentrale.
  \item Die Schl"usselzentrale erzeugt einen zuf"alligen
        Sitzungsschl"ussel $K$ und sendet das Tupel
        $E_A(Bob,R_A,K,T_B),E_B(Bob,K,T_B),R_B$ an Alice.
  \item Alice entschl"usselt und "uberpr"uft, ob $R_A$ mit dem Wert im
        1-ten Schritt "ubereinstimmt. Alice sendet
        $E_B(Alice,K,T_B),E_K(R_B)$ an Bob.
  \item Bob entschl"usselt und erh"alt den Sitzungsschl"ussel $K$. Er
        "uberpr"uft, ob $T_B$ und $R_B$ mit den Werten im 2-ten
        Schritt "ubereinstimmen.
\end{enumerate}

%%
%% Digital Cash
%%

\section{Digital Cash}

Ein ideales System f"ur "`Digital Cash"' soll folgende sechs
Eigenschaften erf"ullen:
\begin{itemize}
  \item {\em Unabh"angigkeit}: das digitale Bargeld kann "uber Netze
        verschickt und gespeichert werden.
  \item {\em Sicherheit}: das digitale Bargeld kann nicht kopiert und
        mehrfach benutzt werden.
  \item {\em Nichzur"uckverfolgbarkeit}: Die Privatsph"are des
        Benutzers ist gewahrt. Keiner kann feststellen, welche
        Transaktionen vom Benutzer vorgenommen wurden.
  \item {\em Off-line Bezahlung}: Bei der Bezahlung mit dem digitalen
        Bargeld ist keine Kommunikation mit der Bank notwendig. Die
        Bezahlung erfolgt off-line.
  \item {\em "Ubertragbarkeit}: Das digitale Bargeld kann an andere
        Benutzer "ubertragen werden.
  \item {\em Teilbarkeit}: Das digitale Bargeld kann in kleinere
        Betr"age aufgeteilt werden.
\end{itemize}
Das erste vorgeschlagene Protokoll, das alle diese Punkte erf"ullte,
ben"otigt 200 MByte Datentransfer f"ur eine Buchung. Mitterweile
existieren Protokolle, die nur 20 KByte Datentransfer brauchen und in
wenigen Sekunden arbeiten.

Im folgenden pr"asentieren wir ein Protokoll, das die ersten vier
Eigenschaft erf"ullt. Wir werden das Protokoll schrittweise ausbauen,
um Sicherheitsl"ucken zu beseitigen.

\subsection{Erste Protokolle}

\subsubsection*{1. Protokoll}
\begin{enumerate} 
  \item Alice stellt 100 anonyme Bankschecks "uber jeweils 1000 DM aus.
  \item Alice steckt jeden Bankscheck zusammen mit Durchschlagspapier
        in einen Umschlag und tr"agt dann alle Umschl"age zur Bank.
  \item Die Bank "offnet 99 Umsch"age und pr"uft, ob der richtige
        Betrag (1000 DM) auf jedem Scheck von Alice eingetragen wurde.
  \item Die Bank unterzeichnet den letzten unge"offneten Umschlag. Die
        Unterschrift der Bank wird durch das Durchschlagspapier auf
        den Scheck "ubertragen.  Die Bank gibt den unge"offneten
        Umschag an Alice zur"uck und belastet ihr Konto mit 1000 DM.
  \item Alice "offnet den Umschag und verwendet den Scheck bei einem
        H"andler.
  \item Der H"andler "uberpr"uft die Unterschrift der Bank und
        akzeptiert den Scheck.
  \item Der H"andler bringt den Scheck zur Bank.
  \item Die Bank "uberpr"uft ihre Unterschrift und zahlt 1000 DM dem
        H"andler aus.
\end{enumerate}
Bei diesem Protokoll kann die Bank nicht feststellen, von wem der
Scheck kommt. Die Wahrscheinlichkeit, da\3 ein Betrug von der Bank
unentdeckt bleibt, betr"agt durch die 1-aus-100-Wahl lediglich
$1/100$. Durch die Einf"uhrung einer Geldstrafe f"ur Betrug lohnt es
sich nicht f"ur Alice zu betr"ugen.

Allerdings gibt es noch eine Sicherheitsl"ucke, da Alice oder der
H"andler den Scheck kopieren und mehrfach verwenden k"onnten. Mit dem
folgenden Protokoll soll diese Sicherheitsl"ucke beseitigt werden.


\subsubsection*{2. Protokoll} 

\begin{enumerate}
  \item Alice stellt 100 anonyme Bankschecks "uber jeweils 1000 DM
        aus. Auf jeden Scheck tr\"agt sie eine Zufallszahl ein, die lang
        genug ist, um die Wahrscheinlichkeit, da\3 eine andere Person
        dieselbe Zufallszahl benutzt, verschwindened klein ist.
  \item Alice steckt jeden Bankscheck zusammen mit Durchschlagspapier
        in einen Umschlag und tr"agt alle Umschl"age zur Bank.
  \item Die Bank "offnet 99 Umsch"age und pr"uft, ob der richtige
        Betrag (1000 DM) und eine jeweils andere Zufallszahl auf jedem
        Scheck von Alice eingetragen wurde.
  \item Die Bank unterzeichnet den letzten unge"offneten Umschlag. Die
        Unterschrift der Bank wird durch das Durchschlagspapier auf
        den Scheck "ubertragen.  Die Bank gibt den unge"offneten
        Umschag an Alice zur"uck und belastet ihr Konto mit 1000 DM.
  \item Alice "offnet den Umschag und verwendet den Scheck bei einem
        H"andler.
  \item Der H"andler "uberpr"uft die Unterschrift der Bank und
        akzeptiert den Scheck.
  \item Der H"andler bringt den Scheck zur Bank.
  \item Die Bank "uberpr"uft ihre Unterschrift und schaut in einer
        Datenbank nach, ob ein Scheck mit derselben Zufallszahl
        eingel"ost wurde. Falls nicht, dann zahlt die Bank dem
        H"andler 1000 DM aus. Die Bank speichert die Zufallszahl in
        ihrer Datenbank ab.
  \item Falls aber ein Scheck mit derselben Zufallszahl eingel"ost
        wurde, dann akzeptiert die Bank den Scheck nicht.
\end{enumerate}
Die Bank kann jetzt zwar einen Betrug feststellen, kann aber den 
Betr"uger nicht identifizieren. Der H"andler ist im Nachteil, da er nicht 
feststellen kann, ob Alice ihm eine Kopie des Schecks aush"andigt. 

\subsubsection*{3. Protokoll}

\begin{enumerate}
  \item Alice stellt 100 anonyme Bankschecks "uber jeweils 1000 DM
        aus. Auf jeden Scheck tr\"agt sie eine Zufallszahl ein, die lang
        genug ist, um die Wahrscheinlichkeit, da\3 eine andere Person
        dieselbe Zufallszahl benutzt, verschwindened klein ist.
  \item Alice steckt jeden Bankscheck zusammen mit Durchschlagspapier
        in einen Umschlag und tr"agt alle Umschl"age zur Bank.
  \item Die Bank "offnet 99 Umsch"age und pr"uft, ob der richtige
        Betrag (1000 DM) und eine jeweils andere Zufallszahl auf jedem
        Scheck von Alice eingetragen wurde.
  \item Die Bank unterzeichnet den letzten unge"offneten Umschlag. Die
        Unterschrift der Bank wird durch das Durchschlagspapier auf
        den Scheck "ubertragen.  Die Bank gibt den unge"offneten
        Umschag an Alice zur"uck und belastet ihr Konto mit 1000 DM.
  \item Alice "offnet den Umschag und verwendet den Scheck bei einem
        H"andler.
  \item Der H"andler "uberpr"uft die Unterschrift der Bank.
  \item Der H"andler fordert Alice auf einen Zufallsstring auf den
        Scheck zu schreiben.
  \item Alice tut dies.
  \item Der H"andler bringt den Scheck zur Bank.
  \item Die Bank "uberpr"uft ihre Unterschrift und schaut in einer
        Datenbank nach, ob ein Scheck mit derselben Zufallszahl
        eingel"ost wurde. Falls nicht, dann zahlt die Bank dem
        H"andler 1000 DM aus. Die Bank speichert die Zufallszahl und den
        Zufallsstring in ihrer Datenbank ab.
  \item Falls aber ein Scheck mit derselben Zufallszahl eingel"ost
        wurde, dann akzeptiert die Bank den Scheck nicht. Sie
        "uberpr"uft den Zufallsstring auf dem Scheck mit dem in der
        Datenbank. Falls sie "ubereinstimmen, wei\3 die Bank, da\3 der
        H"andler den Scheck kopiert hat. Falls sie verschieden sind,
        dann wei\3 die Bank, da\3 die Person, die den Scheck gekauft
        hat, ihn kopiert hat.
\end{enumerate}

Damit ist ein Betrug des H\"andlers, der sich der Bank zu erkennen
gibt, nicht mehr m\"oglich, die anonyme Alice ist aber in der Lage zu
betr\"ugen.
Um Alices Anonymit\"at zu sichern, solange sie sich korrekt verh\"alt,
ihre Identit\"at aber offenzulegen, sobald sie versucht, zu betr\"ugen
bedarf es einiger Hilfsmittel, die jetzt bereitgestellt werden. Weiterhin
ist noch eine elektronische Version der "`Unterschrift mit
Durchschlagspapier"' zu finden.

\subsection{Hilfsprotokolle}

\subsubsection*{Blinde Unterschrift}

Die Bank soll eine Nachricht $m$ von Alice ungesehen
unterschreiben. F"ur die Signatur wird das RSA-Signaturverfahren
benutzt. Die Bank verwendet ihren geheimen Schl"usse~$d$. Der
"offentliche Schl"ussel $(e,n)$ ist allgeimein zug"anglich.

Alice berechnet 
$$ t := m \cdot k^e \pmod{n},$$
wobei $k$ eine Zufallszahl ist.
Die Bank signiert $t$, indem sie $t^d$ berechnet und das Ergebnis
Alice liefert. Alice berechnet
$$ s := t^d \cdot k^{-1} \equiv m^d \cdot (k^e)^d \cdot k^{-1}
   \equiv m^{d} \cdot k^1 \cdot k^{-1} \equiv m^d \pmod{n} $$
und erh"alt die mit dem geheimen Schl"ussel der Bank signierter
Nachricht $m^d$. Die Bank kennt die Nachricht $m$ nicht.

Die blinde Unterschrift entspricht dem Unterschreiben des Umschlags
durch die Bank, wobei die Unterschrift durch das Durchschlagspapier
auf den Scheck "ubertragen wird.


\subsubsection*{Verschleiern einer Nachricht}

Alice w"ahlt eine Zufallszahl $k$ und ver"offentlicht $(f(I,k),k)$,
wobei $f$ eine bekannte Einwegfunktion ist. Jeder kann nach der
Bekanntgabe von $I$ verifizieren, da\3 auch wirklich $f(I,k)$
ver"offentlicht wurde.

Mit dieser Technik l"a"st sich das Problem, die Anonymit"at einer
ehrlichen Alice zu sichern und bei einem Betrug doch die Identit"at
des Betr"ugers zu erhalten:

Alice spaltet eine Bitfolge $I$, die ihre Identit"at vollst"andig
beschreibt (Name, Adresse und alle anderen Informationen, die die Bank
verlangt), in zwei Bitfolgen derart auf, da\3 $I = I_L \oplus I_R$ gilt.
Dann ver\"offentlicht Alice einige dieser "`halbierten Identit\"aten"' in der Form
$\left(f\left(I_{i_L},k_{i_1}\right),k_{i_1}\right)$ und
$\left(f\left(I_{i_R},k_{i_2}\right),k_{i_2}\right)$.
Kann man es erreichen, da\3 bei korrektem Verhalten von jedem Paar nur
eine H\"alfte offengelegt wird, bei einem Betrug jedoch von mindestens
einem Paar beide Teile, so ist diese scheinbar widerspr\"uchliche
Anforderung von Anonymit\"at und Bekanntgabe der Identit\"at
erf\"ullt.


\subsection{Ein vollst\"andiges Protokoll}

\begin{enumerate}
  \item Alice stellt $n$ anonyme Bankschecks "uber einen festen,
        jeweils identischen Betrag aus. Auf jeden Scheck tr\"agt sie
        eine gen\"ugend lange Zufallszahl ein.
        Auf jedem Scheck wird $m$ mal der Identit\"atsstring von
        Alice eingetragen, jeweils auf verschiedene Weisen halbiert 
        und verschleiert (s.o.).
  \item Alle $n$ Schecks werden von Alice f\"ur eine blinde
        Unterschrift der Bank vorbereitet.
  \item Alice "offnet nach Wahl der Bank $n-1$ dieser Schecks;
        die Bank pr\"uft den Betrag, die Zufallszahl und alle (von Alice
        auch offengelegten) Identit\"atsstrings.
  \item Die Bank signiert den letzten Scheck, gibt ihn Alice und belastet
        ihr Konto.
  \item Alice entfernt die Verschleierung und verwendet den Scheck bei einem
        H"andler.
  \item Der H"andler "uberpr"uft die Unterschrift der Bank.
  \item Der H"andler fordert Alice auf von den $m$ Identit\"atsstrings
        jeweils, nach seiner Wahl, die linke oder rechte H\"alfte bekanntzugeben.
  \item Alice tut dies.
  \item Der H"andler bringt den Scheck zur Bank.
  \item Die Bank "uberpr"uft ihre Unterschrift und schaut in einer
        Datenbank nach, ob ein Scheck mit derselben Zufallszahl
        eingel"ost wurde. Falls nicht, dann zahlt die Bank dem
        H"andler den Betrag aus. Die Bank speichert die Zufallszahl in
        ihrer Datenbank ab.
  \item Falls aber ein Scheck mit derselben Zufallszahl eingel"ost
        wurde, dann akzeptiert die Bank den Scheck nicht. Sie
        vergleicht die offengelegten H"alften der Identit"atsstrings 
        mit denen des Schecks in der Datenbank. Falls sie alle "ubereinstimmen, 
        wei\3 die Bank, da\3 der
        H"andler den Scheck kopiert hat. Falls sie verschieden sind,
        dann wei\3 die Bank, da\3 die Person, die den Scheck gekauft
        hat, ihn kopiert hat.
        Da der zweite H\"andler h\"ochst wahrscheinlich eine andere 
        Kombination der halbierten Identit\"aten hat offenlegen lassen,
        kennt die Bank jetzt von einem der Identit\"atsstrings beide 
        H\"alften und damit die Identit\"at des Betr\"ugers.
\end{enumerate}


  